%%=============================================================================
%% METODOLOGIA
%%=============================================================================
\chapter{Metodología}
\label{cap:metodologia}

El desarrollo del proyecto se divide en tres fases: diseño, implementación y validación. Las tres no se ejecutaron de forma estrictamente secuencial, sino como ciclos cortos: cada resultado de validación que contradecía una hipótesis de diseño (por ejemplo, que el \textit{jitter} de ejecución existe en cualquier procesador) dio lugar a un cambio de diseño, a una nueva implementación y a una nueva ronda de pruebas. Este capítulo describe el enfoque general, el objetivo de cada fase, las plataformas utilizadas, los experimentos realizados y los criterios con los que se interpretaron los resultados. El detalle técnico de lo construido se presenta en el Capítulo~\ref{cap:diseno} y los resultados en el Capítulo~\ref{cap:resultados}.

\section{Enfoque de la investigación}

La investigación es aplicada de carácter experimental y cuantitativo. Su pregunta central, si QPP-RNG puede proveer entropía acreditable para criptografía ligera en dispositivos IoT con recursos restringidos, se descompuso en las cuatro preguntas PI1 a PI4 y en las hipótesis H1 y H2 formuladas en la Introducción. La Tabla~\ref{tab:preguntas_experimentos} relaciona cada pregunta con los experimentos que la responden, descritos en la Sección~\ref{sec:experimentos}.

\begin{table}[H]
\centering
\caption{Relación entre preguntas de investigación, hipótesis y experimentos}
\label{tab:preguntas_experimentos}
\small
\begin{tabularx}{\textwidth}{@{}L{0.09\textwidth}Y L{0.12\textwidth}L{0.16\textwidth}@{}}
\toprule
\textbf{Pregunta} & \textbf{Formulación abreviada} & \textbf{Hipótesis} & \textbf{Experimentos} \\
\midrule
PI1 & ¿Reproduce la referencia el comportamiento publicado? & -- & E4 (batería), E6 \\
PI2 & ¿Aporta el \textit{jitter} de ejecución la entropía contabilizada, medida en la fuente? & H1 & E1, E4 \\
PI3 & ¿Qué fuente física la sustituye, cuánto se le acredita y a qué costo? & H2 & E2, E3 \\
PI4 & ¿Se integra con Ascon-AEAD128 en un microcontrolador de 8 bits? & H2 & E5, E6 \\
\bottomrule
\end{tabularx}
\end{table}

Un principio rector de todo el trabajo es que una fuente de entropía no se valida observando únicamente su salida. Como se discutió en la Sección~\ref{sec:evaluacion_estadistica}, superar NIST SP 800-22 o ENT es una evidencia complementaria, pero no suficiente. Un generador determinístico bien diseñado también las supera. Por ello, cada afirmación sobre entropía se apoyó en tres tipos de evidencia: (i) la comparación de capturas independientes a partir del mismo estado inicial, (ii) la estimación de min-entropía de la fuente cruda con NIST SP 800-90B \cite{Turan2018RecommendationGeneration}, condicionada al estado del generador cuando la fuente tiene realimentación, y (iii) las pruebas estadísticas de la salida. La primera es una versión reducida de las pruebas de reinicio de SP 800-90B: basta para detectar una fuente sin ninguna variación entre reinicios, pero no para estimar cuánta entropía hay entre ellos, de modo que el procedimiento formal de la norma, con mil reinicios, queda como trabajo futuro.

\section{Fase de diseño}

La fase de diseño partió de la propuesta de QPP-RNG \cite{Kuang2025QPP-RNG:Randomness,Kuang2025IID-BasedJitter} y de su evaluación multiplataforma \cite{Vrana2025RawEvaluation} y tuvo tres objetivos.

El primero fue establecer una \textbf{implementación de referencia} que sirviera de base. Dado que la publicación describe el mecanismo de forma narrativa y no publica pseudocódigo ni código fuente, cada decisión de implementación que el artículo deja abierta (la forma de la caminata de convergencia, la combinación de los ciclos de sobremuestreo, la evolución de la semilla y los generadores internos) se documentó explícitamente junto al código, para que pudiera contrastarse de forma independiente con el artículo.

El segundo fue definir una \textbf{arquitectura modular} que separara las responsabilidades que la literatura suele presentar juntas, es decir, la medición del tiempo, el generador de referencia, la variante optimizada para IoT, el acondicionamiento criptográfico y el cifrado autenticado. La separación tiene un propósito metodológico. La fuente cruda, que es lo que SP 800-90B caracteriza, debe poder evaluarse sin modificaciones, y cada variante debe poder compararse con la referencia con las mismas herramientas.

El tercero fue definir una \textbf{matriz de evaluación} común para los candidatos del arnés. Las dimensiones de esa matriz se tomaron del Marco teórico (Tabla~\ref{tab:metricas_evaluacion}) y se ampliaron con dos dimensiones que resultaron decisivas: la salud de la fuente de \textit{jitter} en cada plataforma y la reproducibilidad entre reinicios. La Tabla~\ref{tab:matriz_evaluacion} resume la matriz.
 Las herramientas de comparación del computador anfitrión (pruebas estadísticas, de rendimiento, de huella y pruebas diferenciales) cubren a los candidatos registrados en el arnés, es decir, la referencia y su versión acondicionada. La variante IoT, cuya entropía solo existe en las placas, se evaluó con sus propias pruebas unitarias y con las pruebas sobre hardware.
\begin{table}[H]
\centering
\caption{Matriz de evaluación del generador}
\label{tab:matriz_evaluacion}
\begin{tabularx}{\textwidth}{@{}L{0.27\textwidth}L{0.33\textwidth}Y@{}}
\toprule
\textbf{Dimensión} & \textbf{Pregunta} & \textbf{Instrumento} \\
\midrule
Aleatoriedad (IID) & ¿Es la salida independiente e idénticamente distribuida? & NIST SP 800-90B, \texttt{ea\_iid} \\
Min-entropía (no IID) & ¿Cuánta min-entropía puede acreditarse de forma conservadora? & NIST SP 800-90B, \texttt{ea\_non\_iid} \\
Calidad estadística general & ¿Presenta la salida patrones detectables? & NIST SP 800-22 (STS 2.1.2), ENT, pruebas nativas \\
Salud de la fuente & ¿Varía la fuente por causas físicas en el dispositivo real? & Capturas entre reinicios, estimación de la fuente cruda \\
Desempeño & Bytes por segundo, latencia, tiempo al primer byte & \textit{Criterion}, contadores de ciclos \\
Huella & Flash, RAM estática, tamaño por \textit{crate} & \texttt{cargo size}, \texttt{cargo bloat}, \texttt{avr-size} \\
Corrección & ¿Es la implementación determinística dada la misma entrada y libre de pánicos? & Pruebas diferenciales y basadas en propiedades \\
Integración & ¿Produce claves y \textit{nonces} utilizables por Ascon-AEAD128? & Pruebas de integración, vectores KAT y ACVP \\
\bottomrule
\end{tabularx}
\end{table}

\section{Fase de implementación}

La implementación se realizó en Rust (edición 2024, versión mínima 1.97.1), en tres repositorios independientes dentro de la organización QPP-IoT-Lib:

\begin{itemize}
    \item \textbf{QPP-RNG-IoT}: los \textit{crates} del generador (referencia, variante IoT, acondicionamiento, temporizadores y rasgos comunes), el arnés de pruebas y el \textit{firmware} de pruebas para hardware.
    \item \textbf{ASCON-IoT}: una implementación base de Ascon-AEAD128 conforme a NIST SP 800-232 \cite{Turan2025Ascon-BasedFunctions}, desacoplada de cualquier generador.
    \item \textbf{QPP-IoT}: la fachada que combina los dos anteriores y añade las fuentes de entropía específicas de cada placa.
\end{itemize}

Todo el código del camino de generación es \texttt{no\_std} (no depende de la biblioteca estándar ni de memoria dinámica), de modo que el mismo código compila para el computador anfitrión, para Cortex-M y para AVR. La implementación siguió un orden fijo, de modo que cada pieza nueva se probara sobre una base ya validada: primero los rasgos comunes y los temporizadores, luego la referencia, después el arnés de pruebas completo, más adelante el \textit{firmware} de pruebas sobre hardware y la variante IoT y, por último, la integración con ASCON.

\section{Fase de validación}

\subsection{Plataformas}

La Tabla~\ref{tab:plataformas} enumera las plataformas utilizadas. Los cuatro microcontroladores cubren el espectro del IoT restringido: en un extremo, un núcleo AVR de 8 bits sin multiplicador de 64 bits ni divisor; en el intermedio, un Cortex-M4F de 32 bits con generador de números aleatorios en hardware y, en el otro, un ESP32 de doble núcleo Xtensa LX6 a 240 MHz que ejecuta el código desde una memoria Flash externa a través de una caché, la plataforma más difundida en dispositivos IoT con conectividad Wi-Fi. La ESP32 permite comprobar si una caché basta para introducir la variabilidad que falta en los núcleos más simples. Los anfitriones sirven como control y como plataformas en las que la referencia sí dispone de las fuentes de variabilidad que la literatura describe.

\begin{table}[H]
\centering
\caption{Plataformas de evaluación}
\label{tab:plataformas}
\begin{tabularx}{\textwidth}{@{}L{0.25\textwidth}L{0.32\textwidth}Y@{}}
\toprule
\textbf{Plataforma} & \textbf{Procesador / reloj} & \textbf{Uso} \\
\midrule
Arduino Mega 2560 & ATmega2560, AVR de 8 bits, 16 MHz & Pruebas sobre hardware (HIL) \\
Arduino Nano & ATmega328P, AVR de 8 bits, 16 MHz, 2 KB de RAM & Pruebas HIL \\
makerdiary nRF52840 MDK & nRF52840, Cortex-M4F, 64 MHz & Pruebas HIL \\
Placa de desarrollo ESP32 & ESP32-D0WD-V3 (rev.~3.1), Xtensa LX6 de doble núcleo, 240 MHz, 4 MB de Flash & Pruebas HIL \\
Computador Apple Silicon & macOS, ARM64 & Desarrollo, arnés completo, control \\
Raspberry Pi 4 Model B & Cortex-A72 ×4, 0.6--1.5 GHz, Debian 13 & Referencia en un anfitrión rápido \\
Computador de escritorio & Intel Core Ultra 7 265KF, Windows 11 & Referencia en un anfitrión rápido, capturas largas \\
Computador portátil & Intel Core i5-1135G7, WSL2 & Línea base de desempeño de ASCON \\
\bottomrule
\end{tabularx}
\end{table}

Antes de escribir cualquier \textit{firmware}, se respaldó el contenido original de la memoria Flash de las cuatro placas, de modo que cada experimento fuera reversible. El respaldo de la ESP32 se conservó solo de forma local y no se versionó, porque su partición de configuración contenía entradas de red inalámbrica que podrían incluir credenciales.

\subsection{Selección de las fuentes físicas}
\label{sec:seleccion_fuentes}

La fuente física de cada placa se seleccionó con cinco criterios, en este orden: (i) que su variación provenga de un fenómeno físico independiente del reloj que alimenta la CPU, de modo que el modo de falla de la referencia quede excluido por construcción; (ii) que esté disponible en la placa sin hardware adicional; (iii) que el fabricante la documente; (iv) que pueda capturarse de forma cruda desde el \textit{firmware}, para evaluarla con SP 800-90B; y (v) que su tasa de muestras sea compatible con el uso previsto. Con esos criterios, la selección difiere entre placas.

\begin{itemize}
    \item \textbf{AVR.} Los ATmega328P y ATmega2560 no tienen RNG en hardware. El único oscilador interno independiente del cristal de 16 MHz que está disponible sin componentes externos es el oscilador RC de 128 kHz del perro guardián, que puede medirse con el Timer1, alimentado por el cristal. Otras alternativas se descartaron: el ruido de un pin analógico sin conectar depende del montaje y del entorno y está al alcance de un atacante cercano, y el estado de la SRAM al encendido \cite{Holcomb2009SRAM} solo se renueva tras un ciclo completo de alimentación, no tras un reinicio.
    \item \textbf{nRF52840 y ESP32.} Ambos chips incluyen un RNG en hardware documentado por el fabricante como fuente de ruido físico \cite{NordicNRF52840PS,EspressifRNG}, con una tasa mucho mayor que la de un oscilador de baja frecuencia. Por los criterios (iii) y (v), se seleccionó el RNG en estas dos placas.
\end{itemize}

En consecuencia, \emph{la comparación entre placas no compara una misma arquitectura de fuente}. En AVR la fuente es el \textit{jitter} de un oscilador RC, con una acreditación de 1 bit por muestra; en el nRF52840 y en la ESP32 es un RNG en hardware, cuyo diseño interno no es público y al que se acreditan 4 bits por muestra. Lo que sí es común a las cuatro placas es la tubería: pruebas de salud, contabilidad explícita, caminata y acondicionamiento. Las diferencias de rendimiento y de entropía entre placas deben leerse, por tanto, como diferencias entre fuentes y no como una propiedad de la variante IoT. Para disponer, además, de una comparación de la misma arquitectura de fuente en distintas plataformas, en el nRF52840 y en la ESP32 también se caracterizó un oscilador RC de baja frecuencia medido con el cristal (el LFRC y el reloj lento del dominio RTC, respectivamente), análogo al perro guardián de AVR (experimento E2).

\subsection{Representación de las muestras}
\label{sec:representacion}

Los estimadores de SP 800-90B admiten símbolos de hasta 8 bits, mientras que las muestras crudas son más anchas: el periodo del perro guardián, de unos 270\,000 ciclos, ocupa 19 bits, y la palabra del RNG de la ESP32, 32. Por ello la min-entropía de las fuentes se estimó sobre el byte de menor orden de cada muestra, y la de la medición de tiempo de la referencia, sobre un residuo de 8 bits (experimento E4). Esta representación es apropiada por tres razones. Primera, coincide con lo que el diseño consume: la variante IoT aplica las pruebas de salud al byte bajo de cada muestra y la referencia absorbe $\Delta t \bmod 256$, de modo que el valor acreditado corresponde a la misma cantidad sobre la que se configuran las pruebas de salud. Segunda, en una medición de periodo la variación física se concentra en los bits bajos, mientras que los bits altos codifican el valor nominal y su deriva lenta. Tercera, descartar bits solo puede reducir la entropía, nunca aumentarla: el byte bajo es una cota inferior de la información contenida en la muestra completa. En AVR, la variante IoT incorpora la muestra completa a su semilla; en el nRF52840 y en la ESP32, la muestra que usa la variante es el propio byte del RNG, de modo que la representación evaluada y la consumida son idénticas.

El riesgo de esta representación es que la cuantización o el pliegue módulo 256 oculten estructura, por ejemplo, que una deriva lenta de los bits altos aparezca en el byte bajo como una rampa que se pliega y parezca más uniforme de lo que es. Para evitarlo se aplicaron cuatro salvaguardas: (i) además de la estimación por símbolos, SP 800-90B evalúa cada bit por separado ($8\,H_{\text{bitstring}}$) y se reporta el mínimo de ambas, que en varias de las fuentes resultó ser la estimación por bits; (ii) la deriva y la correlación se analizaron sobre la muestra completa (periodo medio, desviación estándar y autocorrelación con desfase 1), y no solo sobre el byte bajo; (iii) se compararon capturas de 20\,000 y de 1\,000\,001 muestras, porque los estimadores de tuplas y de predicción solo detectan la deriva lenta en capturas largas; y (iv) el valor acreditado se fijó claramente por debajo de la estimación, para absorber la estructura que los estimadores pudieran no detectar. En el residuo de la referencia, la representación se eligió para que la parte determinada por la semilla no pase al símbolo evaluado (Sección~\ref{sec:res_anfitriones}); como los residuos quedan casi siempre a pocos \textit{ticks} de cero, el pliegue módulo 256 solo afecta a valores atípicos de la cola de la distribución.

\subsection{Niveles de prueba}

La validación se organizó en cinco niveles de costo creciente:

\begin{enumerate}
    \item \textbf{Pruebas unitarias} de cada \textit{crate}, ejecutadas con \texttt{cargo test}.
    \item \textbf{Pruebas diferenciales y basadas en propiedades}: determinismo de cada configuración frente a un reloj simulado, paridad de comportamiento entre candidatos y \textit{fuzzing} con \textit{proptest} para los anchos de permutación $N \in \{1, 2, 5, 6, 8\}$ en \texttt{xorshift128+} y $N \in \{1, 5, 6\}$ en NEXT\_X48. La variante IoT no forma parte de estas pruebas y se verifica mediante sus propias pruebas unitarias.
    \item \textbf{Pruebas estadísticas de nivel 1}: cinco pruebas nativas en Rust (frecuencia monobit, rachas, $\chi^2$ de uniformidad de bytes con 255 grados de libertad, correlación serial y entropía de Shannon), con nivel de significancia $\alpha = 0.01$, ejecutadas sobre la muestra completa.
    \item \textbf{Pruebas estadísticas de nivel 2}: las herramientas de referencia del NIST (\texttt{ea\_iid}, \texttt{ea\_non\_iid} y \texttt{assess} de STS 2.1.2) y ENT \cite{Walker2008ENT:Program}, sobre muestras de al menos 1\,000\,000 de bytes, el tamaño mínimo recomendado por SP 800-90B.
    \item \textbf{Pruebas sobre hardware} (\textit{hardware-in-the-loop}, HIL): \textit{firmware} dedicado que transmite por UART al anfitrión los valores internos del generador o de la fuente física, para analizarlos con los niveles anteriores.
\end{enumerate}

\subsection{Experimentos}
\label{sec:experimentos}

Se diseñaron seis experimentos. Cada uno responde a una de las preguntas de investigación y genera datos que los demás reutilizan.

\paragraph{E1. Determinismo de la referencia entre reinicios.} En cada microcontrolador se ejecutó la referencia con su configuración por defecto ($N = 5$, \texttt{xorshift128+}), pero con sobremuestreo 1, de modo que cada ciclo de convergencia enviara al anfitrión su par $(n_p, \Delta t)$. Como la salida por defecto es el XOR de cinco valores consecutivos de $n_p \bmod 256$, esta traza determina por completo la salida. El experimento consiste en reiniciar la placa varias veces con la misma semilla y comparar las trazas posición por posición. Si la fuente aporta entropía, las trazas deben diferir desde los primeros ciclos; si no, son idénticas. El computador anfitrión se utilizó como control.

Las condiciones de E1 fueron las siguientes: el mismo binario por placa, con la semilla fijada en tiempo de compilación; reinicios por la línea de reinicio de cada placa, sin cortar la alimentación; alimentación por USB desde el anfitrión, a temperatura ambiente de laboratorio; un solo núcleo en ejecución, sin sistema operativo, sin radio y sin periféricos activos más allá del temporizador y de la UART que transmite la traza; y, en la ESP32, el código ejecutado desde la Flash externa a través de la caché, con el núcleo 1 detenido.

\paragraph{Alcance de la evidencia de determinismo.} Si las trazas coinciden bit a bit, lo que se concluye es que, \emph{en esas condiciones}, $\Delta t$ es una función determinística del estado del programa, es decir, que la perturbación $\varepsilon$ de la Ecuación~\eqref{eq:ciclo_dt} es nula o no alcanza a cambiar un solo valor de la traza. Es una evidencia fuerte, porque una coincidencia exacta en miles de registros es incompatible con cualquier fuente de ruido de magnitud apreciable, y porque tiene una explicación mecánica: el temporizador cuenta el mismo reloj que la CPU y nada externo al programa interviene. No demuestra, en cambio, que toda ejecución de esa plataforma sea determinística bajo cualquier condición de operación. Una aplicación real puede añadir interrupciones de periféricos, acceso directo a memoria, actividad de la radio, un segundo núcleo, ciclos completos de alimentación o condiciones ambientales extremas, y cualquiera de ellos podría introducir variación en $\Delta t$. Esa posibilidad, sin embargo, no rehabilita la referencia como fuente de entropía: una variación que depende de la carga de la aplicación no está caracterizada, desaparece cuando la carga desaparece (por ejemplo, durante el arranque, cuando se generan las claves) y puede estar al alcance de un adversario que controla el tráfico o el entorno. Para acreditarla habría que medirla, condicionada al estado, en cada configuración de despliegue, que es exactamente lo que el diseño original no hace.

\paragraph{E2. Caracterización de fuentes físicas independientes.} En cada microcontrolador se midió, en ciclos del reloj principal, el periodo de un oscilador físicamente independiente de la CPU. En AVR, la interrupción del perro guardián (\textit{watchdog}) se alimenta con un oscilador RC de 128 kHz y marcas de tiempo del Timer1, que cuenta con el cristal de 16 MHz. En el nRF52840, el RTC está se alimenta por el oscilador RC de baja frecuencia (LFRC) y por el contador de ciclos, además del periférico RNG del chip. En la ESP32, el reloj lento del dominio RTC, un oscilador RC de unos 156 kHz, medido en ciclos del cristal de 40 MHz por la unidad de calibración del grupo de temporizadores 0, además del periférico RNG del chip con la fuente de ruido del convertidor analógico-digital (ADC) activada. La razón de seleccionar estas fuentes, y no otras, se expone en la Sección~\ref{sec:seleccion_fuentes}. Primero se tomaron 20\,000 muestras por fuente, para una estimación preliminar y luego 1\,000\,001, para una estimación sin advertencias de tamaño de muestra. La min-entropía $\hat{H}_\infty$ se estimó con \texttt{ea\_non\_iid} sobre el byte de menor orden de cada muestra, con las salvaguardas de la Sección~\ref{sec:representacion}. El valor acreditado $h$ a cada fuente se fijó por debajo de $\hat{H}_\infty$, redondeando hacia abajo y dejando un margen.

\paragraph{E3. Variante IoT sobre hardware.} Con la fuente de E2 y el valor acreditado, se ejecutó la variante \texttt{qpp-rng-iot} en las cuatro placas. Se midieron la divergencia entre reinicios, el rendimiento, las fallas de las pruebas de salud y la huella en Flash, y se aplicaron las pruebas de nivel 1, ENT y SP 800-90B en la salida. Se capturaron primero muestras cortas (10\,000 bytes por AVR y 50\,000 en el nRF52840) y luego muestras de 1\,000\,000 de bytes por placa, para la batería completa. La ESP32, evaluada al final, pasó directamente a la muestra de 1\,000\,000 de bytes, tras una verificación corta de la divergencia entre reinicios.

\paragraph{E4. Entropía por ciclo de referencia en anfitriones rápidos.} Se repitió la traza de E1 en una Raspberry Pi 4 y en un computador con Windows 11, con 5\,000\,005 y 2\,000\,000 ciclos, respectivamente. Como el valor $n_p$ de cada ciclo ya está determinado por la semilla, la entrada nueva que recibe la semilla en cada ciclo es $\Delta t$ \emph{condicionado} a $n_p$. Por ello, la min-entropía fresca por ciclo, $H_\infty(\delta \mid s)$ en la notación de la Ecuación~\eqref{eq:entropia_fresca}, se estimó con \texttt{ea\_non\_iid} sobre el residuo $\Delta t - \operatorname{moda}(\Delta t \mid n_p)$ y con un estimador directo de la Ecuación~\eqref{eq:min_entropia_condicional}, y se comparó con la entropía que el diseño contabiliza. El procedimiento y sus supuestos se detallan junto con los resultados (Sección~\ref{sec:res_anfitriones}). Además, se generaron muestras de 1\,000\,000 de bytes de los cuatro candidatos de la librería, para aplicarles la batería completa.

\paragraph{E5. Validación de Ascon-AEAD128.} La implementación base se verificó contra el conjunto completo de 1\,089 vectores de respuesta conocida (KAT), contra una validación local de extremo a extremo con el servidor ACVP del NIST, y de forma diferencial contra la implementación de referencia en C (\texttt{ascon-c}), incluidas pruebas negativas de autenticación.

\paragraph{E6. Desempeño y huella.} Se midieron el rendimiento y la latencia de los candidatos en el anfitrión con \textit{Criterion}, los ciclos por byte en los microcontroladores con el Timer1 o el contador de ciclos, el tamaño del \textit{firmware}, el tamaño de la pila completa (fuente, generador, SHA-256 y Ascon) en AVR, y el tiempo de cifrado de Ascon-AEAD128 en el anfitrión. Por último, la pila completa se ejecutó de extremo a extremo en las cuatro placas: generación de clave y \textit{nonce} desde la fuente física, cifrado, descifrado y rechazo de mensajes alterados, con cada texto cifrado verificado contra \texttt{ascon-c} en el anfitrión, y con el tiempo de cada etapa y el consumo máximo de pila medidos en el dispositivo.

\subsection{Criterios de interpretación}

Para que los resultados no dependan de lecturas \textit{ad hoc}, se fijaron de antemano los siguientes criterios:

\begin{itemize}
    \item \textbf{Divergencia entre reinicios.} Dos capturas desde el mismo estado inicial se consideran independientes cuando el número de coincidencias posicionales es compatible con el azar. Para bytes uniformes, se esperan $n/256$ coincidencias en $n$ posiciones; para $n_p \sim \text{Geom}(1/120)$, se esperan alrededor de 83 en 20\,000. Una coincidencia del 100\,\% indica una fuente sin entropía.
    \item \textbf{Pruebas con múltiples subpruebas.} Con $\alpha = 0.01$, una fuente ideal falla, en promedio, el 1\,\% de las combinaciones de prueba y flujo. Por ejemplo, en SP 800-22, con 8 flujos de 1\,Mbit, se esperan unas 14.5 fallas entre 1\,452 combinaciones. Se considera anómalo un número de fallas claramente superior a ese valor, o fallas que se repiten de forma sistemática en una misma prueba.
    \item \textbf{Min-entropía acreditada.} El valor $h$ que se acredita a una fuente es menor o igual que la estimación $\hat{H}_\infty$ de \texttt{ea\_non\_iid} sobre al menos 1\,000\,000 de muestras, redondeado hacia abajo. Las estimaciones de SP 800-90B sobre la salida acondicionada no se usan para acreditar la entropía, porque una función hash borra precisamente los patrones que esos estimadores buscan.
    \item \textbf{Tamaño de muestra.} Los resultados obtenidos sobre muestras menores que las recomendadas se reportan como pruebas de humo y no como validación, y se contrastan con controles del mismo tamaño tomados de \texttt{/dev/urandom}.
\end{itemize}

\subsection{Herramientas}

La Tabla~\ref{tab:herramientas} resume las herramientas utilizadas. Las herramientas del NIST se compilaron desde sus fuentes oficiales en macOS, en la Raspberry Pi y en WSL; la equivalencia entre las instalaciones se comprobó reproduciendo, en WSL, fila por fila, un informe de SP 800-22 obtenido en macOS.

\begin{table}[H]
\centering
\caption{Herramientas de construcción, captura y análisis}
\label{tab:herramientas}
\begin{tabularx}{\textwidth}{@{}L{0.32\textwidth}Y@{}}
\toprule
\textbf{Herramienta} & \textbf{Uso} \\
\midrule
Rust 1.97.1--1.99.0 (estable), \textit{nightly} para \texttt{avr-none} & Compilación del anfitrión y de los microcontroladores (\texttt{-Z build-std=core} en AVR) \\
avr-gcc 14.3.0, avrdude 8.3 & Rutinas en C para AVR, programación de las placas Arduino \\
probe-rs 0.32.0 & Programación del nRF52840 mediante CMSIS-DAP \\
espup 0.19.0 (Rust 1.97.0 para Xtensa), \texttt{esp-hal} 1.2.2, espflash 4.7.0 & Compilación y programación de la ESP32 \\
NIST SP 800-90B (\texttt{ea\_iid}, \texttt{ea\_non\_iid}) \cite{NISTEntropyAssessment} & Estimación de min-entropía \\
NIST STS 2.1.2 (\texttt{assess}) \cite{Rukhin2010AApplications} & Batería SP 800-22 \\
ENT (Fourmilab) & Estadísticas por byte \\
\textit{Criterion}, \textit{proptest} & Medición de desempeño, pruebas basadas en propiedades \\
\texttt{cargo size}, \texttt{cargo bloat} & Huella en memoria \\
NIST ACVP-Server \cite{NISTACVPServer}, \texttt{ascon-c} \cite{AsconC} & Validación de Ascon-AEAD128 \\
\bottomrule
\end{tabularx}
\end{table}

\section{Amenazas a la validez}

Se identificaron las siguientes limitaciones, que condicionan la generalización de los resultados:

\begin{itemize}
    \item \textbf{Validez externa.} Se evaluó un ejemplar de cada placa. La entropía de un oscilador RC varía entre ejemplares, temperaturas y tensiones de alimentación, de modo que el valor acreditado debe reevaluarse para cada modelo de dispositivo que se despliegue.
    \item \textbf{Tamaño de muestra.} Las muestras cortas de E3 están por debajo de lo recomendado por SP 800-90B. Por ello se repitieron con muestras de 1\,000\,000 de bytes, y sus resultados se reportan por separado.
    \item \textbf{Fidelidad de la referencia.} La referencia es una interpretación documentada de un artículo sin código publicado. Las diferencias con la implementación original de los autores no pueden descartarse; solo pueden hacerse explícitas.
    \item \textbf{Entorno de medición.} Las mediciones de desempeño de Ascon se realizaron en WSL2, un entorno virtualizado sin control de frecuencia, y no son extrapolables a los microcontroladores. En la ESP32, el código se ejecuta desde la Flash a través de una caché, y el \textit{firmware} de extremo a extremo mide cada operación una sola vez, con la caché fría, de modo que sus tiempos de cifrado no son comparables con los de las demás placas.
    \item \textbf{Tipo de reinicio.} Los reinicios se realizaron mediante la línea de reinicio de cada placa (la señal DTR del puerto serie en los Arduino, una señal de interrupción serie en el nRF52840 y el pin EN, mediante la señal RTS, en la ESP32). No se evaluaron ciclos completos de alimentación, en los que el estado inicial de los osciladores y de la memoria podría ser distinto.
    \item \textbf{Pruebas de reinicio reducidas.} Las comparaciones entre reinicios usaron entre 2 y 13 reinicios por placa, no los mil que prescribe el procedimiento de reinicio de SP 800-90B. Son suficientes para establecer que la referencia no varía en absoluto y que la variante IoT sí varía desde el primer byte, pero no para estimar la min-entropía entre reinicios.
    \item \textbf{Fuentes distintas por placa.} Como se explicó en la Sección~\ref{sec:seleccion_fuentes}, la fuente de AVR es un oscilador RC y la del nRF52840 y la ESP32 es un RNG en hardware. Los resultados de la variante IoT no son, por tanto, una comparación de la misma fuente en cuatro arquitecturas.
    \item \textbf{Validación formal.} Las validaciones con ACVP y con las herramientas del NIST son locales y reproducibles, pero no constituyen una certificación CAVP ni una validación de la fuente de entropía ante el NIST.
    \item \textbf{Origen de la variabilidad en AVR.} No se aisló el ruido térmico de la deriva por temperatura o tensión (autocorrelación de 0.93 a 0.97 en el millón de muestras). Un atacante con acceso físico podría influir en esa deriva; la contabilidad de 1 bit por muestra es conservadora frente a las estimaciones de \texttt{ea\_non\_iid} (1.62 y 2.24 bits).
    \item \textbf{Ritmo de lectura del RNG de la ESP32.} El RNG se evaluó leyéndolo una vez cada 1.65 ms, mientras que en operación se lee a unas 1\,540 muestras por segundo y, al arrancar, se le realizan ocho lecturas seguidas; las pruebas de salud no saltaron, pero la evaluación no replica exactamente ese patrón.
\end{itemize}
